% !TeX program = pdflatex
% Compile twice with pdfLaTeX. No external bibliography is required.
% Revision: 8 October 2026.
% The compression theorem is quoted from the pinned OpenAI formal source.
% Its proof is not reproduced or independently machine-checked here.
\documentclass[11pt,a4paper]{article}
\usepackage[T1]{fontenc}
\usepackage[utf8]{inputenc}
\usepackage{lmodern}
\usepackage[margin=26mm]{geometry}
\usepackage{amsmath,amssymb,amsthm,mathtools}
\usepackage{microtype}
\usepackage[hidelinks]{hyperref}
\usepackage{bookmark}
\setlength{\emergencystretch}{2em}
\numberwithin{equation}{section}
\newtheorem{theorem}{Theorem}
\newtheorem{proposition}[theorem]{Proposition}
\newtheorem{lemma}[theorem]{Lemma}
\newtheorem{external}[theorem]{External theorem}
\newcommand{\Fq}{\mathbb F_q}
\newcommand{\HH}{\mathrm H}
\newcommand{\EE}{\mathbb E}
\newcommand{\cC}{\mathcal C}
\newcommand{\cS}{\mathcal S}
\newcommand{\eps}{\varepsilon}
\newcommand{\src}[1]{\texttt{\detokenize{#1}}}
\newcommand{\oairoot}{https://github.com/openai/math/blob/adc7f1241b42e322a6451854ab7e4b4c146bf78a/lean/OAI/Combinatorics/SharpRamsey/}
\newcommand{\source}[2]{\href{\oairoot#1}{\src{#2}}}
\hypersetup{pdftitle={Common independent sets in graphs with different forbidden cliques}}
\title{Common independent sets in $K_{s_1}$-free, $\ldots$, $K_{s_k}$-free graphs}
\author{}
\date{}

\begin{document}
\maketitle
\vspace{-2em}

\begin{abstract}
Using the single-colour compression theorem in OpenAI's formal source,
we prove that, for every fixed integer $k\ge1$, fixed integers
$s_1,\ldots,s_k\ge3$, and every $\eps>0$,
$r(s_1,\ldots,s_k,t)
\ge c_{\boldsymbol{s},\eps}t^{S+1}/(\log t)^{S+\eps}$,
where $\boldsymbol{s}=(s_1,\ldots,s_k)$ and
$S=\sum_{i=1}^k(s_i-2)$.
Together with the upper bound of He--Wigderson, this determines the
logarithmic exponent. We use Brada\v{c}'s ordered flag construction.
A rectangle-peeling argument removes the occupancy assumption from the
quoted compression theorem, and projection-closed covering families
allow the resulting reduction to be iterated across the colours.
Triangle-free colours are handled by the independent-set count of
Andrades--Campos--Morris, applied uniformly over the retained lengths.
The external theorem, its specialization, and the additional arguments
are stated separately.
\end{abstract}

\section{The construction}

Fix an integer $k\ge1$ and integers $s_1,\ldots,s_k\ge3$, and write
$\boldsymbol{s}=(s_1,\ldots,s_k)$ and $S=\sum_{i=1}^k(s_i-2)$.
The Ramsey number $r(s_1,\ldots,s_k,t)$ is the least $N$
such that every $(k+1)$-colouring of $E(K_N)$ contains a
$K_{s_i}$ in colour $i$ for some $1\le i\le k$, or a $K_t$
in colour $k+1$.

\begin{theorem}\label{thm:main}
For every fixed integer $k\ge1$, fixed integers $s_1,\ldots,s_k\ge3$,
and every $\eps>0$, there is a constant
$c_{\boldsymbol{s},\eps}>0$ such that, for all sufficiently large integers $t$,
\[
 r(s_1,\ldots,s_k,t)
 \ge c_{\boldsymbol{s},\eps}\frac{t^{S+1}}{(\log t)^{S+\eps}}.
\]
Consequently,
$r(s_1,\ldots,s_k,t)=t^{S+1}/(\log t)^{S+o(1)}$ as $t\to\infty$.
\end{theorem}

A parameter $s_i=2$ may be deleted. If all the small parameters
are $2$, the Ramsey number equals $t$.

The upper bound used here is the unconditional estimate in
He--Wigderson~\cite[Appendix~B]{HW}.
For the lower bound, our external compression input is
\source{Marking/SelectedStep.lean}{eventually_selected_step}
in the formal source~\cite{OAI}, specialized to each fixed vector
dimension $s_i\ge4$. The colours with $s_i=3$ use
\cite[Theorem~1.3]{ACM}, as detailed in Section~\ref{sec:triangles}.
We state this input in Section~\ref{sec:compression}, together with
its source notation and hypotheses, before deriving the covering
estimate used in the proof. We do not reprove the external compression
argument.

All logarithms are natural. For a nonnegative integer $n$, write
$[n]=\{1,\ldots,n\}$, with $[0]=\emptyset$.
Constants may depend on fixed parameters, but not on $q$.

For the one-colour arguments, fix an integer $d\ge4$.
Let $q$ be prime, and recall that $PG(d-1,q)$ is the set of
one-dimensional subspaces of $\Fq^d$. For $a,b\in PG(d-1,q)$,
we write $a\perp b$ if their nonzero representative vectors have
zero dot product. This condition does not depend on the representatives.
Let $D=D^*(d-1,q)$ be the digraph introduced by Brada\v{c}
\cite[Section~2.5]{Bradac}. Its vertex set is
$V(D)=\{(a,b)\in PG(d-1,q)^2:a\perp b\}$,
and there is an arc from $(a,b)$ to $(a',b')$ if and only if
$a\perp b'$ and $a'\not\perp b$.

For any positive integer $N$ and any map
$\phi:[N]\to V(D)$, let $G_\phi$ be the graph on $[N]$
with edge set
$E(G_\phi)=\{uv:1\le u<v\le N,\ \phi(u)\phi(v)\in A(D)\}$.
An ordered tuple $(x_1,\ldots,x_m)$ of vertices of $D$,
where repetitions are allowed, is called \emph{forward independent}
if there is no arc from $x_i$ to $x_j$ whenever $i<j$.

For each subspace $U\le\Fq^d$, let $U^\perp$ denote its
orthogonal complement, and define
$R(U)=\{(a,b)\in V(D):a\subseteq U^\perp,\ b\subseteq U\}$.
We call these sets \emph{orthogonal rectangles}, using the terminology
of the compression framework in~\cite{OAI}.

\begin{lemma}\label{lem:construction}
The following properties hold.
\begin{enumerate}
\item[(i)] For every positive integer $N$ and every map
$\phi:[N]\to V(D)$, the graph $G_\phi$ is $K_d$-free.

\item[(ii)] For $u_1<\cdots<u_m$ in $[N]$, the set
$\{u_1,\ldots,u_m\}$ is independent in $G_\phi$ if and only if
$(\phi(u_1),\ldots,\phi(u_m))$ is forward independent.
Moreover, writing $x_i=(a_i,b_i)$, a tuple
$(x_1,\ldots,x_m)$ is forward independent if and only if the
following holds: for every $i<j$, if $a_i\perp b_j$, then
$a_j\perp b_i$.

\item[(iii)] Every order-preserving subsequence of a forward
independent tuple is forward independent.

\item[(iv)] Writing $v=|V(D)|$, we have, for every subspace
$U\le\Fq^d$,
\begin{equation}\label{eq:geometry}
 q^{2d-3}\le v\le4q^{2d-3},\qquad
 |R(U)|\le4q^{d-2},\qquad
 |\{U:U\le\Fq^d\}|\le q^{d^2}.
\end{equation}
\end{enumerate}
\end{lemma}

\begin{proof}
By \cite[Lemma~2.11, with parameter $d-1$]{Bradac}, $D$ contains
no transitive tournament on $d$ vertices.
If $u_1<\cdots<u_d$ formed a clique in $G_\phi$, then
$\phi(u_1),\ldots,\phi(u_d)$ would be distinct, since $D$ has no
loops, and would form such a tournament. This proves (i).
Properties (ii) and (iii) follow directly from the definitions.

For (iv), counting projective points gives
$v=(q^d-1)(q^{d-1}-1)/(q-1)^2$.
If $\dim U=e\in\{0,\ldots,d\}$, then
$|R(U)|=(q^e-1)(q^{d-e}-1)/(q-1)^2$.
These formulas imply the first two bounds.
Finally, every subspace of $\Fq^d$ is the row space of some
$d\times d$ matrix, giving at most $q^{d^2}$ subspaces.
\end{proof}

\section{Products and the entropy bound}\label{sec:products}

Fix $\eps>0$ and an integer $L=L(\boldsymbol{s},\eps)\ge5$, to be chosen at
the end of the proof, and set
\begin{equation}\label{eq:scale}
 T=\lfloor q(\log q)^{1+4/L}\rfloor,\qquad
 N=\lfloor q^{S+1}\log q\rfloor.
\end{equation}
After relabelling the first $k$ colours, assume that $s_i\ge4$
for $1\le i\le a$ and $s_i=3$ for $a<i\le k$, where $0\le a\le k$.
Put $b=k-a$. For $1\le i\le a$, write
$D_i=D^*(s_i-1,q)$ and $v_i=|V(D_i)|$.
Our goal is to construct $K_{s_i}$-free graphs $G_i$ on $[N]$
for $1\le i\le k$ such that $\alpha(G_1\cup\cdots\cup G_k)<T$.
For $i\le a$ we use maps $\phi_i:[N]\to V(D_i)$ and
$G_i=G_{\phi_i}$; the remaining graphs are specified in
Section~\ref{sec:triangles}. In the one-colour arguments below,
$D$ continues to denote $D^*(d-1,q)$.

For an integer $0\le m\le T$, a \emph{product of length $m$}
is a set $C=C_1\times\cdots\times C_m$, where each
$C_i$ is a nonempty subset of $V(D)$.
The product of length zero consists of the single empty tuple.
We say that $C$ \emph{covers} $(x_1,\ldots,x_m)$ if
$x_i\in C_i$ for every $i\in[m]$.
A family $\cC$ of products covers a tuple if some member of
$\cC$ of the same length covers it. All covering families below
are finite and nonempty, and may contain products of different
lengths, all at most $T$.

For $I\subseteq[m]$, write $I=\{i_1,\ldots,i_s\}$ with
$i_1<\cdots<i_s$. The \emph{coordinate projection} of $C$ onto
$I$ is $\pi_I(C)=C_{i_1}\times\cdots\times C_{i_s}$.
When $I=\emptyset$, this is the product of length zero.
A family $\cC$ is \emph{projection-closed} if it contains every
coordinate projection of each of its members.
For a tuple $x=(x_1,\ldots,x_m)$, write
$x_I=(x_{i_1},\ldots,x_{i_s})$.

Fix sufficiently large constants $P,A$, depending on $d$; the requirements
on these constants are specified below. For the multicolour argument,
choose them to satisfy these requirements simultaneously for the
finitely many dimensions $d=s_i\ge4$.
An admissible covering pair $(\cC,\Delta)$ consists of a covering
family $\cC$ and a common parameter $\Delta\ge0$ such that every
factor of every member satisfies $|C_s|\le Pq^{d-1}e^\Delta$.
Its \emph{covering complexity} is
\begin{equation}\label{eq:domains}
 B(\cC,\Delta)=\log|\cC|+Aq(\log q)\Delta.
\end{equation}
We write $B$ when the pair is understood.
Thus a statement about a family with complexity $B$ includes the
existence of its stated common domain parameter; complexity is
attached to the pair, not to the family alone.

For a random tuple $X$, let $\HH(X)$ denote its Shannon entropy,
computed using natural logarithms. We use the chain rule,
subadditivity, and the bound by the logarithm of the support size;
see~\cite[Chapter~2]{CT}.
The trivial estimate is $\HH(X)\le m\log v=(2d-3)m\log q+O(m)$.
The following proposition replaces the leading term by
$(d-1)m\log q$, at the additional cost of the covering complexity.
It is useful when that complexity is small.

\begin{proposition}[One-colour entropy bound]\label{prop:entropy}
Let $(\cC,\Delta)$ be an admissible covering pair, and let $X$
be a random forward independent tuple of length $1\le m\le T$
whose support is covered by $\cC$.
For all sufficiently large $q$,
\begin{equation}\label{eq:entropy-upper}
 \HH(X)\le(d-1)m\log q+Am+B(\cC,\Delta).
\end{equation}
This holds for every probability distribution of $X$;
projection-closure and rectangle-occupancy assumptions are not required.
\end{proposition}

\begin{proof}
By \cite[Claim~2.13, with parameter $d-1$]{Bradac}, there is a
constant $b=b(d)>0$ such that the tree of forward independent
tuples can be marked so that every node has at most $bq^{d-1}$ marked
children and every path from the root has at most $bq\log q$
unmarked vertices. Fix an integer $\kappa\ge\max\{P,b,1\}$,
and put $w=\lceil bq\log q\rceil$.

Fix a product $C=C_1\times\cdots\times C_m$ in $\cC$.
Restrict the tree to branches allowed by these coordinate sets.
The marking properties are preserved, and total branching is at
most $\kappa q^{d-1}e^\Delta$.
If $m\ge w$, \cite[Lemma~2.7]{Bradac} bounds the number of
permitted forward independent tuples by
$(2\kappa q^{d-1})^m e^{w\Delta}$.
If $m<w$, the trivial bound $(\kappa q^{d-1}e^\Delta)^m$ is no larger.
Summing over products of length $m$ and bounding entropy by the
logarithm of the support size gives
$\HH(X)\le(d-1)m\log q+m\log(2\kappa)+w\Delta+\log|\cC|$.
Taking $A\ge\max\{\log(2\kappa),b+1\}$ proves the result.
\end{proof}

\section{The external compression input}\label{sec:compression}

Throughout this section we use the following revision of~\cite{OAI}:
\begin{center}
\src{adc7f1241b42e322a6451854ab7e4b4c146bf78a}.
\end{center}
We first state the selected-step theorem in its
source variables, retaining its occupancy and excess-cost
hypotheses. We then verify the specialization needed here.
The latter verification, including passage from an auxiliary
probability law to a finite cover, is part of this paper's argument.

\subsection{The quoted theorem}

Put $V=\Fq^d$ and $\mathcal P=PG(V)\times PG(V^*)$.
Let $\mathcal F\subseteq\mathcal P$ be the set of incident
point--dual-point pairs $(p,[\lambda])$, where $\lambda$ vanishes
on $p$. A tuple $F_s=(p_s,[\lambda_s])$ is
\emph{scan-consistent} if, whenever $s<t$ and
$\lambda_s(p_t)=0$, we have $\lambda_t(p_s)=0$.
This is \src{ScanConsistent} after the source conversion
\src{toScan}. For $W\le V$, write
$\mathcal R(W)=\{(p,[\lambda])\in\mathcal F:
 p\subseteq W,\ W\subseteq\ker\lambda\}$.

We need one numerical quantity from the source theorem.
For a tuple $f\in\mathcal F^n$, let $E_q(f)\subseteq[n]$ be the
set of positions marked expensive in either of the source's two
canonical scans. Precisely, this is
\source{Marking/MarkingTwoScans.lean}{bothExpensive}
applied to \src{twoMask} at parameter $q$; the second scan is
reversed and dualized, and its positions are returned to the
original indexing.
Given a finite probability space $(\Omega,p)$, a tuple-valued
map $F:\Omega\to\mathcal F^n$, a finite-valued message $u$,
and a real number $J$, set $J_0=(d-1)\log q+\log64$ and
\begin{equation}\label{eq:source-excess-definition}
 \mathfrak e(F,u,J)
 =\HH(u)+2n\log(3(d+1))+nJ_0-\HH(F)
   +(J-J_0)\EE_p|E_q(F)|.
\end{equation}
Formula~\eqref{eq:source-excess-definition} is the specialization of
\source{Entropy/MarkingTwoEntropy.lean}{excessCost}
to vector dimension $d$.
The term $2n\log(3(d+1))$ follows from the mask-cardinality identity
\source{Marking/ExcessCost.lean}{twoMask_log_card}.
We do not need the internal scan construction beyond this definition
and the following quoted estimates.

The source declarations
\source{Marking/ExcessCost.lean}{twoExpensive_linear}
and
\source{Marking/ExcessCost.lean}{excessCost_linear},
with their dimension parameter equal to $d-3$, give the following.
For $\log q\ge1$, the sum of the two source expensive-position
allowances is at most $64(d+1)^3q\log q$.
Moreover, if $P\ge64$, $\Delta,\Lambda,\tau\ge0$,
$J=(d-1)\log q+\log P+\Delta$,
$\HH(u)\le\Lambda$, and
$\HH(F)\ge(d-1)n\log q-\tau n$, then
\begin{equation}\label{eq:source-excess-bound}
 \mathfrak e(F,u,J)
 \le\Lambda+n(\tau+2\log(3(d+1))+\log P)
       +64(d+1)^3q(\log q)\Delta.
\end{equation}
Here the constant-gap parameter in \src{excessCost_linear} is
$\log(P/64)$.

\begin{external}[OpenAI's selected-step theorem in dimension $d$]
\label{thm:oai-step}
Let $P_0=\src{stepPrefactor}(d-1)$ be the domain prefactor, depending only on $d$,
in~\cite{OAI}. Fix $0<\eta\le1$ and positive constants
$c_-,c_+,C_0$, and put $\beta=\eta/10^7$.
For all sufficiently large primes $q$, the following holds,
uniformly over the data below.

Let $(\Omega,p)$ be a nonempty finite probability space,
let $F:\Omega\to\mathcal F^n$ have scan-consistent values on
its support, and let $u:\Omega\to\Gamma$ be a message with a
finite alphabet $\Gamma$. Suppose that nonempty domains
$S(\gamma,s)\subseteq\mathcal P$ satisfy
$F_s(\omega)\in S(u(\omega),s)$ whenever $p(\omega)>0$,
and $\log|S(\gamma,s)|\le J$ for every $\gamma,s$.
Let $\ell\ge1$ be an integer, $D_0$ a real number, and $R_0$
an even positive integer, such that
\begin{equation}\label{eq:source-admissible}
 \begin{gathered}
 (\log q)^\beta\le D_0\le(\log q)^{1-\eta/2},\\
 (\log q)^\beta\le R_0\le2(\log q)^\beta+2,
 \end{gathered}
\end{equation}
and
\begin{equation}\label{eq:source-length}
 c_-q(\log q)^{1+\eta}\le\ell
 \le c_+q(\log q)^{1+\eta},\qquad
 d(d+2)\ell+64(d+1)^3q\log q\le n.
\end{equation}
Assume also that
\begin{equation}\label{eq:source-budget}
 \mathfrak e(F,u,J)
 \le\frac{C_0\ell D_0(\log q)^{-\beta}}{d(d+2)},
\end{equation}
and that, for every $\omega$ with $p(\omega)>0$ and every
subspace $W\le V$,
\begin{equation}\label{eq:source-occupancy}
 |\{s\in[n]:F_s(\omega)\in\mathcal R(W)\}|
 \le C_0\log q.
\end{equation}
Then there are an integer $m'$ with $\ell/64\le m'\le n$,
a finite probability space $(\Xi,\mu)$, maps
$\rho:\Xi\to\Omega$ and $M:\Xi\to\Gamma'$ with
$\Gamma'$ finite, increasing injections
$\iota_z:[m']\to[n]$, and domains
$\mathcal T(\gamma,s)\subseteq\mathcal P$ such that
\begin{enumerate}
\item[(i)] For every $z$ in the support of $\mu$ and every
$s\in[m']$, we have
$F_{\iota_z(s)}(\rho(z))\in \mathcal T(M(z),s)$.
\item[(ii)] For every $\omega\in\Omega$,
$\mu(\rho^{-1}(\omega))\le4d(d+2)p(\omega)$.
\item[(iii)] The new message and domains satisfy
\begin{equation}\label{eq:source-output}
 \begin{gathered}
 \HH_\mu(M)\le nD_0(\log q)^{-\eta/3},\\
 |\mathcal T(\gamma,s)|\le
 P_0q^{d-1}\exp\bigl(16D_0(\log q)^{6\beta}\bigr).
 \end{gathered}
\end{equation}
\end{enumerate}
The domains depend only on the new message and the new coordinate,
not on the old message $u$ as an additional argument.
The threshold for $q$ depends only on $d,\eta,c_-,c_+,C_0$.
\end{external}

This is
\source{Marking/SelectedStep.lean}{eventually_selected_step}
with $i=d-4$. Its class count is $d(d+2)$, and its domination factor is
$4d(d+2)$.
The packing condition in \eqref{eq:source-length} is sufficient
for the original packing condition by \src{twoExpensive_linear}.
The inequalities in \eqref{eq:source-admissible}, including the
parity of $R_0$, are precisely
\source{Parameters/SourceScales.lean}{Admissible}.
The source quantity \src{scaleKstar} is
$D_0(\log q)^{6\beta}$.
The output above spells out
\source{Execution/ContextOutput.lean}{ContextOutput};
in particular, its \src{chosen} field is order-preserving,
its \src{dominated} field is the bound in (ii), and its
\src{domains} field has no old-message argument.
No stream-domination assumption is imposed on $p$ in this formal
declaration.

\subsection{The finite-cover consequence}

Choose the constants from Section~\ref{sec:products} once and for
all so that $P\ge\max\{64,P_0\}$ and
$A\ge\max\{10^4,64(d+1)^3,2\log(3(d+1))\}+2\log P$, as well as the requirements in the proof of
Proposition~\ref{prop:entropy}.
For a tuple, \emph{occupancy} counts positions rather than distinct
labels: its occupancy in $R(U)$ is
$|\{s:x_s\in R(U)\}|$.

\begin{lemma}[Dispersed compression]\label{lem:dispersed}
Fix $\delta\in(0,1]$, $K\ge1$, and an integer
$1\le j\le L-4$.
For all sufficiently large primes $q$, let $\Omega$ be a
nonempty finite set with the uniform distribution, and let
$X:\Omega\to V(D)^m$ take forward independent values.
Suppose that an admissible covering pair $(\cC,\Delta)$ covers
$X$ and satisfies
\begin{equation}\label{eq:dispersed-input}
 \delta T\le m\le T,\qquad
 B(\cC,\Delta)\le Km(\log q)^{j/L},\qquad
 \HH(X)\ge(d-1)m\log q-Km(\log q)^{j/L}.
\end{equation}
Suppose additionally that every $X(\omega)$ has at most
$K\log q$ positions in each orthogonal rectangle.
Then there are $\Omega'\subseteq\Omega$, an integer $m'$, and
an admissible covering pair $(\mathcal M,\Delta')$, with every
member of $\mathcal M$ of length $m'$, such that
\begin{equation}\label{eq:dispersed-output}
 |\Omega'|\ge\frac{|\Omega|}{8d(d+2)},\qquad
 m'\ge\frac{m}{256d(d+2)},\qquad
 B(\mathcal M,\Delta')\le m(\log q)^{(j-1)/L}.
\end{equation}
For each $\omega\in\Omega'$, some order-preserving subsequence
of $X(\omega)$ of length $m'$ is covered by $\mathcal M$.
The threshold is uniform over the input laws and covers after
$d,L,\delta,K,j$ are fixed.
\end{lemma}

\begin{proof}
We verify the hypotheses of Theorem~\ref{thm:oai-step} and then
convert its output into a cover.
Put $r=j/L$, $\eta=4/L$, and $\beta=\eta/10^7$.
Identify a label $(a,b)\in V(D)$ with
$\iota(a,b)=(b,[\lambda_a])\in\mathcal F$, where
$\lambda_a(x)=\widetilde a\cdot x$ for any nonzero
representative $\widetilde a$ of $a$.
This is a bijection. Under it, forward independence becomes
scan-consistency, and $R(U)$ becomes $\mathcal R(U)$.
Thus entropy, domain sizes, order, and rectangle occupancies are
preserved. This also agrees with the source rectangle convention
in \source{Construction/InitialTools.lean}{rectangle_pair_iff}.

For each $\omega$, choose a covering product of length $m$
by a fixed rule, and use its index as $u(\omega)$.
Then $\HH(u)\le\log|\cC|$.
The images of its factors supply the domains in the external
theorem. Take $F=\iota\circ X$ coordinatewise,
$n=m$, and $J=(d-1)\log q+\log P+\Delta$.
By \eqref{eq:source-excess-bound},
\begin{equation}\label{eq:actual-excess}
 \begin{aligned}
 \mathfrak e(F,u,J)
 &\le\log|\cC|+Km(\log q)^r
      +(2\log(3(d+1))+\log P)m+64(d+1)^3q(\log q)\Delta\\
 &\le(2K+A)m(\log q)^r.
 \end{aligned}
\end{equation}
The last inequality uses the bound on $B(\cC,\Delta)$ and our
fixed choice of $A$.

Choose $\ell=\lfloor m/[2d(d+2)]\rfloor$,
$D_0=(\log q)^{r+\beta}$, and
$R_0=2\lceil(\log q)^\beta/2\rceil$.
Since $r\le1-\eta$, condition \eqref{eq:source-admissible}
holds for sufficiently large $q$.
Also $m/[4d(d+2)]\le\ell\le m/[2d(d+2)]$ and
$d(d+2)\ell+64(d+1)^3q\log q\le m$, because
$m\ge\delta T\ge(\delta/2)q(\log q)^{1+\eta}$.
Consequently \eqref{eq:source-length} holds with
$c_-=\delta/[8d(d+2)]$ and $c_+=1$.
Take $C_0=4d^2(d+2)^2(2K+A)$.
Then \eqref{eq:actual-excess} implies
\eqref{eq:source-budget}, and the occupancy hypothesis holds
because $C_0\ge K$.

Theorem~\ref{thm:oai-step} now supplies a finite auxiliary law,
a source map $\rho$, a message $M$, and a common output length
$m'\ge\ell/64\ge m/[256d(d+2)]$, with
\begin{equation}\label{eq:raw}
 \begin{gathered}
 \mu(\rho^{-1}(\omega))\le4d(d+2)/|\Omega|,\\
 \HH_\mu(M)\le h:=m(\log q)^{r+\beta-\eta/3},\\
 |\mathcal T(\gamma,s)|\le Pq^{d-1}
             \exp\bigl(16(\log q)^{r+7\beta}\bigr).
 \end{gathered}
\end{equation}

Let $p_M$ be the distribution of $M$.
Retain the auxiliary outcomes for which
$-\log p_M(M)\le2h+2$.
Since the expectation of $-\log p_M(M)$ is $\HH_\mu(M)$,
Markov's inequality retains probability at least $1/2$.
There are at most $e^{2h+2}$ retained message values, since each
has probability at least $e^{-(2h+2)}$.
By \eqref{eq:raw}, the image $\Omega'$ of the retained support
under $\rho$ has at least $|\Omega|/[8d(d+2)]$ elements.

For each $\omega\in\Omega'$, fix one retained witness over
$\omega$. Its increasing injection selects the required
subsequence. For each retained message, intersect its domains
with $\mathcal F$, pull them back under $\iota$, and take their
product. Discard messages not occurring in the retained support.
The resulting factors are nonempty, their sizes do not increase,
and these products form a fixed-length family $\mathcal M$.
Since the domains depend only on the new message, this family
has at most $e^{2h+2}$ members, with no additional factor for the
old covering product.

Set $\Delta'=16(\log q)^{r+7\beta}$.
Using $m\ge(\delta/2)q(\log q)^{1+\eta}$, we obtain
\begin{align*}
 B(\mathcal M,\Delta')
 &\le2+m(\log q)^r\left(
       2(\log q)^{\beta-\eta/3}
       +\frac{32A}{\delta}(\log q)^{7\beta-\eta}\right)\\
 &\le m(\log q)^{r-1/L}.
\end{align*}
Indeed, $1/L=\eta/4<\eta/3-\beta$ and
$1/L<\eta-7\beta$, so both terms in parentheses are
$o((\log q)^{-1/L})$.
The additive constant is absorbed since $m\to\infty$ and
$r-1/L\ge0$.
All applications of the external theorem use fixed constants;
its uniform threshold therefore gives the asserted uniformity.
\end{proof}

\section{Removing the occupancy hypothesis}\label{sec:reduction}

We now remove the extra hypothesis in Lemma~\ref{lem:dispersed}.
If many positions lie in a short list of rectangles, their labels
already have a small cover. Otherwise we retain a dispersed
subsequence and use Proposition~\ref{prop:entropy} to control its
entropy loss before applying the external compression input.

\begin{lemma}[One-colour reduction]\label{lem:reduction}
There is a constant $c=c(d)\in(0,1)$ with the following
property. Fix $\delta\in(0,1]$, $K\ge1$, and an integer
$1\le j\le L-4$.
For all sufficiently large primes $q$, let $\Omega$ be a
nonempty finite set with the uniform distribution, and let
$X:\Omega\to V(D)^m$ take forward independent values.
Suppose that $X$ is covered by an admissible pair $(\cC,\Delta)$,
where $\cC$ is projection-closed, and that
\begin{equation}\label{eq:reduction-input}
 \delta T\le m\le T,\qquad
 B(\cC,\Delta)\le Km(\log q)^{j/L},\qquad
 \HH(X)\ge(d-1)m\log q-Km(\log q)^{j/L}.
\end{equation}
Then there are $\Omega'\subseteq\Omega$, an integer $m'$, and
an admissible pair $(\cC',\Delta')$ with $\cC'$ projection-closed,
such that
\begin{equation}\label{eq:reduction-output}
 |\Omega'|\ge c|\Omega|,\qquad
 cm\le m'\le m,\qquad
 B(\cC',\Delta')\le(d^2+4)m(\log q)^{(j-1)/L}.
\end{equation}
Each $X(\omega)$ with $\omega\in\Omega'$ has an order-preserving
subsequence of the common length $m'$ covered by $\cC'$.
The selected positions may depend on $\omega$.
The threshold depends only on $d,L,\delta,K,j$ and the fixed
constants $P,A$.
\end{lemma}

\begin{proof}
Write $B=B(\cC,\Delta)$ and put $h=\lceil\log q\rceil$.
For each input, repeatedly choose a rectangle containing at
least $h$ remaining positions, record that rectangle, and remove
all positions currently lying in it. Use a fixed ordering of
subspaces to break ties. There are at most $\lfloor m/h\rfloor$
recorded rectangles. Padding their lists with an empty symbol,
Lemma~\ref{lem:construction}(iv) gives
\begin{equation}\label{eq:peel}
 \begin{aligned}
 \log(\text{number of possible lists})
 &\le\lfloor m/h\rfloor\log(q^{d^2}+1)\le(d^2+1)m,\\
 \left|\bigcup\text{recorded rectangles}\right|
 &\le4q^{d-2}m/h\le4q^{d-1}(\log q)^{4/L}.
 \end{aligned}
\end{equation}
After termination, fewer than $h\le2\log q$ remaining positions
lie in any rectangle.

Put $m_0=\lfloor m/2\rfloor$, so $m_0\ge m/4$ for large $q$.
For each input, at least one of the removed and unremoved parts
has at least $m_0$ positions. Select such a part by a fixed rule.
One of the two choices occurs on a subset $\Omega_0\subseteq\Omega$
of size at least $|\Omega|/2$.
For each input in $\Omega_0$, retain the first $m_0$ positions
of its chosen part.

\paragraph{The removed part.}
Suppose the chosen part is the removed one.
For each recorded list arising on $\Omega_0$, let $S$ be the
union of its rectangles, and include the products $S^s$ for
$0\le s\le m_0$ in $\cC'$.
The set $S$ is nonempty since it contains the retained labels.
This family is projection-closed and covers each retained tuple.
Take $\Delta'=(4/L)\log\log q$.
By \eqref{eq:peel} and $P\ge4$, its domains are admissible, and
\[
 B(\cC',\Delta')
 \le(d^2+1)m+\log(m_0+1)
       +\frac{4A}{L}q(\log q)\log\log q
 \le(d^2+3)m
\]
for sufficiently large $q$.
Here $m\ge\delta T$ and
$q(\log q)\log\log q=o(T)$, with $L$ fixed.
Thus \eqref{eq:reduction-output} holds with
$\Omega'=\Omega_0$ and $m'=m_0$, since $(j-1)/L\ge0$.

\paragraph{The unremoved part.}
Suppose instead that the chosen part is the unremoved one.
We first check its entropy.
Put $r=j/L$, let $Z$ be the indicator of $\Omega_0$, and write
$p_0=\Pr(Z=1)\ge1/2$.
Proposition~\ref{prop:entropy} applies to every conditional law
of the full tuple, with the same cover. Hence
\[
 \HH(X)\le\log2+p_0\HH(X\mid Z=1)
 +(1-p_0)((d-1)m\log q+Am+B),
\]
with the last conditional contribution omitted when $p_0=1$.
Using \eqref{eq:reduction-input}, we obtain
\begin{equation}\label{eq:conditioned}
 \HH(X\mid Z=1)
 \ge(d-1)m\log q-O_{A,K}\bigl(m(\log q)^r\bigr).
\end{equation}
Thus the entropy loss under this conditioning is controlled by
an upper bound on the discarded part, not by a general assertion
that conditioning preserves entropy.

Now work on $\Omega_0$ with its uniform distribution.
Let $I\subseteq[m]$ be the input-dependent set of retained
$m_0$ positions, and put $Y=X_I$ and
$W=X_{[m]\setminus I}$, in their original orders.
The full tuple is determined by $(I,Y,W)$.
Both $Y$ and $W$ are forward independent, and both are covered
by the unchanged family $\cC$ because it is projection-closed.
Consequently,
\[
 \HH(W)\le(d-1)(m-m_0)\log q+A(m-m_0)+B,
 \qquad \HH(I)\le\log\binom m{m_0}\le m\log2.
\]
Together with \eqref{eq:conditioned}, subadditivity gives
\begin{equation}\label{eq:subtuple-entropy}
 \begin{aligned}
 \HH(Y)
 &\ge\HH(X\mid Z=1)-\HH(I)-\HH(W)\\
 &\ge(d-1)m_0\log q-O_{A,K}\bigl(m_0(\log q)^r\bigr).
 \end{aligned}
\end{equation}
The covering complexity is still $B\le4Km_0(\log q)^r$,
and $m_0\ge\delta T/4$.
Every rectangle contains at most $2\log q$ positions of $Y$.
Thus Lemma~\ref{lem:dispersed} applies with lower length
proportion $\delta/4$ and a sufficiently large fixed input
constant, at least $2$, absorbing the constants in
\eqref{eq:subtuple-entropy}.

It gives a set $\Omega'\subseteq\Omega_0$, a common length
$m'$, and a fixed-length admissible cover $(\mathcal M,\Delta')$
with
$|\Omega'|\ge|\Omega_0|/[8d(d+2)]$,
$m'\ge m_0/[256d(d+2)]$, and
$B(\mathcal M,\Delta')\le m_0(\log q)^{(j-1)/L}$.
Let $\cC'$ contain all coordinate projections of all members
of $\mathcal M$, with the same parameter $\Delta'$.
There are at most $2^{m'}$ projections per product, so
\[
 \begin{aligned}
 B(\cC',\Delta')
 &\le B(\mathcal M,\Delta')+m'\log2\\
 &\le2m_0(\log q)^{(j-1)/L}
 \le m(\log q)^{(j-1)/L}.
 \end{aligned}
\]
Here $m'\le m_0$ and $(j-1)/L\ge0$.
The chosen subsequences of $Y$ are also order-preserving
subsequences of the original tuples $X$.
Finally,
$|\Omega'|\ge|\Omega|/[16d(d+2)]$ and $m'\ge m/[1024d(d+2)]$.
Thus $c=1/[1024d(d+2)]$ works in both cases.
All input constants for Lemma~\ref{lem:dispersed} depend only
on the fixed parameters, giving the claimed uniformity.
\end{proof}

\section{The triangle-free colours}\label{sec:triangles}

The compression theorem above requires $d\ge4$.
For the colours with $s_i=3$, we use a triangle-free graph with
few independent sets. The same graph must control every length
that occurs after taking subsequences.

For a graph $H$, let $i_m(H)$ denote the number of its independent
sets of size $m$.

\begin{lemma}[Triangle-free colours]\label{lem:triangles}
There is an absolute constant $C_\triangle\ge1$ such that, for all
sufficiently large primes $q$, there is a triangle-free graph $H$
on $[N]$ satisfying
\begin{equation}\label{eq:triangle-count}
 \frac{i_m(H)}{\binom Nm}\le\left(\frac{C_\triangle}{q}\right)^m
 \qquad\text{for every integer }u\le m\le T,
 \qquad u=\left\lfloor\frac{q\log q}{2}\right\rfloor.
\end{equation}
\end{lemma}

\begin{proof}
By \cite[Theorem~1.3]{ACM}, for all sufficiently large integers $u$
and every integer $n\ge u^2/\log u$, there is a triangle-free graph
on $n$ vertices with at most
$e^{C_0u}\binom{pn}{u}$ independent $u$-sets, where
$p=(\log u)/u$ and $C_0$ is an absolute constant.
Apply this with $n=N$ and $u=\lfloor q(\log q)/2\rfloor$.
The size condition holds because $S\ge1$ and
$u^2/\log u=(1/4+o(1))q^2\log q\le N$.
Also $p=O(1/q)$, so the resulting graph $H$ satisfies
\begin{equation}\label{eq:triangle-initial}
 u!\,i_u(H)\le e^{C_0u}(pN)^u\le(C_1N/q)^u
\end{equation}
for an absolute constant $C_1$.

Fix $m\ge u$. If $i_m(H)=0$, there is nothing to prove.
Otherwise, let $Y=(Y_1,\ldots,Y_m)$ be uniform over all ordered
independent $m$-tuples with distinct entries. Thus
$\HH(Y)=\log(m!\,i_m(H))$.
The chain rule and the fact that conditioning reduces entropy give
\[
 \binom{m-1}{u-1}\HH(Y)
 \le\sum_{J\in\binom{[m]}u}\HH(Y_J)
 \le\binom mu\log(u!\,i_u(H)).
\]
Indeed, expand each $\HH(Y_J)$ in increasing coordinate order and
condition each term additionally on all earlier coordinates;
each coordinate then occurs $\binom{m-1}{u-1}$ times.
Consequently, \eqref{eq:triangle-initial} implies
\begin{equation}\label{eq:triangle-larger}
 m!\,i_m(H)\le(u!\,i_u(H))^{m/u}\le(C_1N/q)^m.
\end{equation}
For all $m\le T$, we have $m\le N/2$ for sufficiently large $q$,
since $T=o(N)$. Writing $(N)_m=N(N-1)\cdots(N-m+1)$, it follows that
\[
 \frac{i_m(H)}{\binom Nm}
 =\frac{m!\,i_m(H)}{(N)_m}
 \le\left(\frac{2C_1}{q}\right)^m.
\]
Taking $C_\triangle\ge\max\{1,2C_1\}$ proves the lemma.
\end{proof}

Fix one such graph $H$. For a permutation $\pi$ of $[N]$, let
$H_\pi$ be the graph in which $uv$ is an edge if and only if
$\pi(u)\pi(v)\in E(H)$.
For any fixed $m$-set $U\subseteq[N]$ and independent uniform
permutations $\pi_1,\ldots,\pi_b$, Lemma~\ref{lem:triangles} gives
\begin{equation}\label{eq:triangle-permutations}
 \Pr\bigl(U\text{ is independent in every }H_{\pi_j}\bigr)
 =\left(\frac{i_m(H)}{\binom Nm}\right)^b
 \le\left(\frac{C_\triangle}{q}\right)^{bm}
 \qquad(u\le m\le T).
\end{equation}
When $b=0$, this probability is $1$.
For each fixed $\delta>0$, all lengths $\delta T\le m\le T$
satisfy $m\ge u$ for sufficiently large $q$.
Thus the same graph $H$ applies throughout the induction below.

\section{The \texorpdfstring{$k$}{k}-colour induction}
\label{sec:induction}

\begin{proof}[Proof of Theorem~\ref{thm:main}]
Let $\cS$ be the set of all inputs
$(\phi_1,\ldots,\phi_a,\pi_1,\ldots,\pi_b)$, where
$\phi_i:[N]\to V(D_i)$ for $1\le i\le a$, and each $\pi_j$
is a permutation of $[N]$.
Thus $|\cS|=(N!)^b\prod_{i=1}^a v_i^N$.
Use common constants $P,A$ for the finitely many dimensions
$d=s_i\ge4$, and fix
\[
 K=2A\max_{1\le i\le k}(s_i-2)
     +\max_{1\le i\le k}(s_i^2+4),\qquad
 c=\min_{1\le i\le k}\frac1{1024s_i(s_i+2)}.
\]
This $c$ is a common lower bound for the retained proportions in
Lemma~\ref{lem:reduction} for every dimension that is used.
Empty sums and tuples have value and entropy zero, respectively;
empty products have value one.

We prove the following assertion for every fixed
$\delta\in(0,1]$ and every vector of integers
$(j_1,\ldots,j_a)\in\{0,\ldots,L-4\}^a$:
for all sufficiently large $q$, the following data do not exist.
There is a set $\Omega\subseteq\cS$ with
$|\Omega|\ge\delta|\cS|$, a single integer
$\delta T\le m\le T$, and, for each
$\omega=(\phi_1,\ldots,\phi_a,\pi_1,\ldots,\pi_b)\in\Omega$,
an increasing list $u_1(\omega)<\cdots<u_m(\omega)$ in $[N]$, such that
$X_i(\omega)=(\phi_i(u_1(\omega)),\ldots,
\phi_i(u_m(\omega)))$ is forward independent for every $1\le i\le a$,
and $\{u_1(\omega),\ldots,u_m(\omega)\}$ is independent in
every $H_{\pi_j}$, $1\le j\le b$.
Each $X_i$ is covered by an admissible pair for dimension $d=s_i$,
$(\cC_i,\Delta_i)$, with $\cC_i$ projection-closed and
\begin{equation}\label{eq:bounds}
 B_i:=B(\cC_i,\Delta_i)
 \le KT(\log q)^{j_i/L}\qquad(1\le i\le a).
\end{equation}
The positions may depend on the whole input, but the list
is common to all colours, and its length $m$ is the same for
every input. We induct on $j_1+\cdots+j_a$.
The constant $K$ stays fixed, while the assertion is made for
every fixed $\delta$.

Equip $\Omega$ with its uniform distribution.
For a prescribed $a$-tuple $x$ of $m$-tuples, first choose its
common increasing set of positions. In the uniform distribution
on $\cS$, the probability of the prescribed labels at these
positions is $\prod_{i=1}^a v_i^{-m}$.
Independently, the probability that these positions are independent
in all the triangle-free colours is at most
$(C_\triangle/q)^{bm}$ by \eqref{eq:triangle-permutations};
here $m\ge u$ for sufficiently large $q$, since $\delta$ is fixed.
Summing over the common position sets and then restricting to
$\Omega$ gives
\begin{equation}\label{eq:mixed-atom}
 \Pr_\Omega\bigl((X_1,\ldots,X_a)=x\bigr)
 \le\delta^{-1}\binom Nm
       \left(\frac{C_\triangle}{q}\right)^{bm}
       \prod_{i=1}^a v_i^{-m}.
\end{equation}
There is no extra ordering factor. Bounding entropy below by
minus the logarithm of the largest atom gives
\begin{align}
 \HH(X_1,\ldots,X_a)
 &\ge m\sum_{i=1}^a\log v_i+bm\log q-\log\binom Nm
       -bm\log C_\triangle+\log\delta\notag\\
 &\ge\left(\sum_{i=1}^a(s_i-1)\right)m\log q
       +\frac4L m\log\log q-O_{\boldsymbol{s},\delta}(m).
 \label{eq:joint}
\end{align}
Indeed, $v_i\ge q^{2s_i-3}$,
$m\ge(\delta/2)q(\log q)^{1+4/L}$, and
$N\le q^{S+1}\log q$ imply
$\log(eN/m)\le S\log q-(4/L)\log\log q+\log(2e/\delta)$.
Now use $\binom Nm\le(eN/m)^m$ and
\[
 \sum_{i=1}^a(2s_i-3)+b-S=\sum_{i=1}^a(s_i-1).
\]
This argument allows arbitrary dependence between the selected
tuples and the chosen positions. Independence is used only in
the full experiment on $\cS$, not after restricting to $\Omega$.

\paragraph{The base case.}
If all $j_i$ are zero, including the case $a=0$,
Proposition~\ref{prop:entropy} and subadditivity imply
\[
 \begin{aligned}
 \HH(X_1,\ldots,X_a)
 &\le\sum_{i=1}^a\HH(X_i)\\
 &\le\left(\sum_{i=1}^a(s_i-1)\right)m\log q+aAm+\sum_{i=1}^aB_i\\
 &=\left(\sum_{i=1}^a(s_i-1)\right)m\log q
    +O_{\boldsymbol{s},\delta}(m),
 \end{aligned}
\]
where $K$ is fixed and $T\le m/\delta$.
This contradicts \eqref{eq:joint} for sufficiently large $q$.

\paragraph{The induction step.}
Otherwise, choose $h\in[a]$ such that
$j_h=\max_i j_i\ge1$.
Applying Proposition~\ref{prop:entropy} to all other colours
with index at most $a$ gives
\begin{equation}\label{eq:marginal}
 \begin{aligned}
 \HH(X_h)
 &\ge\HH(X_1,\ldots,X_a)-\sum_{\substack{1\le i\le a\\i\ne h}}\HH(X_i)\\
 &\ge(s_h-1)m\log q-\sum_{\substack{1\le i\le a\\i\ne h}}B_i
       -O_{\boldsymbol{s},\delta}(m).
 \end{aligned}
\end{equation}
For $a=1$, the sums over other colours are empty.
Since $j_i\le j_h$ and $T\le m/\delta$, we have
$B_h\le(K/\delta)m(\log q)^{j_h/L}$ and
$\sum_{i\ne h}B_i\le((a-1)K/\delta)m(\log q)^{j_h/L}$,
where the sum is over $1\le i\le a$.
Thus the hypotheses of Lemma~\ref{lem:reduction}, with $d=s_h$,
hold for $X_h$, with a sufficiently large fixed input constant
depending on $\boldsymbol{s},\delta,K,A$.

Apply that lemma and retain the same selected subsequence in
every other colour.
For each input, composing the new increasing selection with
its old list gives a common increasing list of original
vertices of a fixed length $m'$.
Every untreated covering family already contains the required projections,
so we keep $(\cC_i,\Delta_i)$ unchanged for $1\le i\le a$, $i\ne h$.
Independence in each triangle-free colour is preserved by taking a subset.
The new complexities satisfy
\begin{equation}\label{eq:descent}
 \begin{gathered}
 B_h'\le(s_h^2+4)m(\log q)^{(j_h-1)/L}
       \le KT(\log q)^{(j_h-1)/L},\\
 B_i'=B_i\qquad(1\le i\le a,\ i\ne h).
 \end{gathered}
\end{equation}
Also $|\Omega'|\ge c\delta|\cS|$ and $m'\ge c\delta T$.
The induction hypothesis at
$(j_1,\ldots,j_h-1,\ldots,j_a)$ and $\delta'=c\delta$
therefore gives a contradiction. It is applied with the uniform
law on $\Omega'$, to which \eqref{eq:mixed-atom} applies afresh.
The input constant in the reduction may increase, but the
output coefficient $s_h^2+4$, and hence the induction constant $K$,
does not.

All thresholds used in this step depend only on fixed numerical
data, not on the actual inputs or covers.
For a fixed index vector and $\delta$, take the maximum of the
reduction threshold and the threshold in the smaller-index
induction hypothesis.
This completes the induction; its depth is at most $a(L-4)$,
independently of $q$.

\paragraph{Starting the induction.}
Suppose every input in $\cS$ had a common independent
$T$-set in all $k$ associated graphs. Choose one by a fixed rule
and list its vertices increasingly. Set $\Omega=\cS$, $m=T$,
and, for every $1\le i\le a$,
$\cC_i=\{V(D_i)^r:0\le r\le T\}$ with
$\Delta_i=(s_i-2)\log q$.
These families are projection-closed and admissible, since
$v_i\le4q^{2s_i-3}\le Pq^{s_i-1}e^{(s_i-2)\log q}$.
Moreover,
\[
 B_i=\log(T+1)+A(s_i-2)q(\log q)^2
 \le\bigl(2A(s_i-2)+1\bigr)T(\log q)^{(L-4)/L}
 \le KT(\log q)^{(L-4)/L}
\]
for sufficiently large $q$.
This contradicts the assertion with $\delta=1$ and
all $j_i=L-4$; for $a=0$, it contradicts the base case directly.
Hence some input yields graphs
$G_i=G_{\phi_i}$ for $1\le i\le a$ and
$G_{a+j}=H_{\pi_j}$ for $1\le j\le b$, satisfying
$\alpha(G_1\cup\cdots\cup G_k)<T$.
Each $G_i$ is $K_{s_i}$-free, by Lemma~\ref{lem:construction}(i)
for $i\le a$ and the triangle-freeness of $H$ otherwise.

Colour $E(G_i)\setminus\bigcup_{h<i}E(G_h)$ with colour $i$
for $1\le i\le k$, and all remaining edges with colour $k+1$.
There is no $K_{s_i}$ in colour $i$ for any $1\le i\le k$ and
no $K_T$ in the last colour, so $r(s_1,\ldots,s_k,T)>N$.

\paragraph{Passing to arbitrary $t$.}
For sufficiently large $t$, Bertrand's postulate gives a prime
$q\in[x,2x]$, where $x=t/[4(\log t)^{1+4/L}]$.
Then $T\le t$, $\log q\sim\log t$, and \eqref{eq:scale} yields
\[
 r(s_1,\ldots,s_k,t)
 \ge c_{\boldsymbol{s},L}\frac{t^{S+1}}
 {(\log t)^{S+4(S+1)/L}}.
\]
Choose an integer $L\ge\max\{5,4(S+1)/\eps\}$.
This proves the asserted lower bound.
The unconditional upper bound in
He--Wigderson~\cite[Appendix~B]{HW}, with
$S=\sum_{i=1}^k(s_i-2)$, is
$r(s_1,\ldots,s_k,t)\le C_{\boldsymbol{s}}t^{S+1}/(\log t)^S$.
Together with the lower bound for every fixed $\eps>0$, this
proves the claimed logarithmic exponent.
\end{proof}

\begin{thebibliography}{9}
\small\raggedright
\setlength{\itemsep}{4pt}

\bibitem{Bradac}
D. Brada\v{c}, \emph{Off-diagonal Ramsey numbers},
\href{https://arxiv.org/abs/2605.28793v3}{arXiv:2605.28793v3},
16 June 2026.
Sections~2.3 and~2.5; Lemmas~2.7 and~2.11, and Claim~2.13.

\bibitem{CT}
T. M. Cover and J. A. Thomas,
\emph{Elements of Information Theory}, Wiley, 1991.
Chapter~2.

\bibitem{HW}
X. He and Y. Wigderson,
\emph{Multicolor Ramsey numbers via pseudorandom graphs},
Electronic Journal of Combinatorics \textbf{27} (2020), no.~1,
Paper No.~1.32.
\href{https://doi.org/10.37236/9071}{doi:10.37236/9071};
\href{https://arxiv.org/abs/1910.06287v3}{arXiv:1910.06287v3}.
The upper bound used here is proved in Appendix~B and does not
require the pseudorandom-graph existence hypothesis of the lower bound.

\bibitem{OAI}
OpenAI, \emph{SharpLogRamsey formal source},
\href{https://github.com/openai/math}{\src{openai/math}},
revision
\href{https://github.com/openai/math/tree/adc7f1241b42e322a6451854ab7e4b4c146bf78a/lean/OAI/Combinatorics/SharpRamsey}{\src{adc7f1241b42e322a6451854ab7e4b4c146bf78a}}.
The external results used are
\source{Marking/SelectedStep.lean}{eventually_selected_step},
\source{Marking/ExcessCost.lean}{excessCost_linear}, and
\source{Marking/ExcessCost.lean}{twoExpensive_linear}.
Their auxiliary definitions and the output structure are linked
in Section~\ref{sec:compression}.
Source statements consulted on 8 October 2026.

\bibitem{ACM}
B. Andrades, M. Campos, and R. Morris,
\emph{On the multicolour Ramsey numbers $R(3,3,k)$},
\href{https://arxiv.org/abs/2609.28455v1}{arXiv:2609.28455v1},
23 September 2026.
The independent-set count used in Section~\ref{sec:triangles}
is Theorem~1.3.

\end{thebibliography}
\end{document}
